\indent ಅವರಿಂದ ಹೊಲವು ಮಾರಲ್ಪಡುವುದು\par

\indent (ಮಾರು + ಅಲ್ಪಡು + ಉ + ವದು)\par

\indent ಮೊದಲನೆಯ ಶ್ರೇಣಿಯಲ್ಲಿರುವ ವಾಕ್ಯಗಳಲ್ಲಿ ಕರ್ತೃಪದವು ಪ್ರಥಮಾ ವಿಭಕ್ತ್ಯಂತವಾಗಿರುವುದು ಮತ್ತು ಕರ್ಮಪದವು ದ್ವಿತೀಯಾ ವಿಭಕ್ತ್ಯಂತವಾಗಿರುವುದು. ಅಲ್ಲದೆ ಈ ವಾಕ್ಯಗಳಲ್ಲಿ ಕ್ರಿಯಾಪದದ ಲಿಂಗ ವಚನಗಳು ಕರ್ತೃಪದದ ಲಿಂಗ ವಚನಗಳಂತೆಯೇ ಇರುವುವು. ಉದಾಹರಣೆಗೆ, “ರಾಮನು” ಎಂಬ ಕರ್ತೃಪದವು ಪುಲ್ಲಿಂಗ ಏಕವಚನವಾಗಿರುವುದರಿಂದ ಕ್ರಿಯಾಪದವೂ ಪುಲ್ಲಿಂಗ ಏಕವಚನ ರೂಪವನ್ನು ಹೊಂದಿ “ಮಾರುತ್ತಾನೆ” ಎಂದಾಗಿರುವುದು. ಈ ಲಕ್ಷಣಗಳುಳ್ಳ ಪ್ರಯೋಗಗಳನ್ನು ಕರ್ತರಿ ಪ್ರಯೋಗಗಳೆನ್ನುವರು.\par

\indent ಎರಡನೆಯ ಶ್ರೇಣಿಯಲ್ಲಿರುವ ವಾಕ್ಯಗಳಲ್ಲಿ ಕರ್ತೃಪದವು ತೃತೀಯಾ ವಿಭಕ್ತ್ಯಂತವಾಗಿರುವುದು ಮತ್ತು ಕರ್ಮಪದವು ಪ್ರಥಮಾ ವಿಭಕ್ತ್ಯಂತವಾಗಿರುವುದು. ಅಲ್ಲದೆ, ಕರ್ಮಪದದ ಲಿಂಗ ವಚನಗಳಂತೆ ಕ್ರಿಯಾಪದವೂ ಇರುವುದು. ಅಲ್ಲದೆ, ವಿಶೇಷವೇನೆಂದರೆ, ಕ್ರಿಯಾಪದಗಳಲ್ಲಿ ಧಾತುವಿನ ಮುಂದೆ “ಅಲ್ಪಡು” ಎಂಬ ಪ್ರತ್ಯಯ ಸೇರಿರುವುದು. ಈ ಲಕ್ಷಣಗಳುಳ್ಳ ಪ್ರಯೋಗಗಳನ್ನು ಕರ್ಮಣಿ ಪ್ರಯೋಗಗಳೆನ್ನುವರು.\par

\vspace{3pt plus 2pt}\noindent\hspace*{1.2em}\kbvlabel{ಉದಾ}\par\vspace{1pt}

\vspace{2.5pt plus 3pt minus 0.5pt}\noindent\hspace*{1.2em}\kbvnum{1.}\hspace{0.55em}ರೈತನು ಭೂಮಿಯನ್ನು ಉಳುತ್ತಾನೆ. {[}ಕರ್ತರಿ ಪ್ರಯೋಗ{]} ರೈತನಿಂದ ಭೂಮಿಯು ಉಳಲ್ಪಡುತ್ತದೆ. {[}ಕರ್ಮಣಿ ಪ್ರಯೋಗ{]}\par

\vspace{2.5pt plus 3pt minus 0.5pt}\noindent\hspace*{1.2em}\kbvnum{2.}\hspace{0.55em}ಬೆಕ್ಕು ಇಲಿಯನ್ನು ಕೊಂದಿತು. {[}ಕರ್ತರಿ ಪ್ರಯೋಗ{]} ಬೆಕ್ಕಿನಿಂದ ಇಲಿಯು ಕೊಲ್ಲಲ್ಪಟ್ಟಿತು. {[}ಕರ್ಮಣಿ ಪ್ರಯೋಗ{]}\par

\vspace{2.5pt plus 3pt minus 0.5pt}\noindent\hspace*{1.2em}\kbvnum{3.}\hspace{0.55em}ಅಗಸರು ಬಟ್ಟೆಗಳನ್ನು ಒಗೆಯುವರು. {[}ಕರ್ತರಿ ಪ್ರಯೋಗ{]} ಅಗಸರಿಂದ ಬಟ್ಟೆಗಳು ಒಗೆಯಲ್ಪಡುವುವು. {[}ಕರ್ಮಣಿ ಪ್ರಯೋಗ{]}\par

\kbvsechead{ನೆನಪಿಡಿ :}

\vspace{2.5pt plus 3pt minus 0.5pt}\noindent\hspace*{1.2em}\kbvnum{(1)}\hspace{0.55em}ಸಕರ್ಮಕ ಧಾತುವಿನ ಮುಂದೆ “ಅಲ್ಪಡು” ಪ್ರತ್ಯಯ ಬಂದ ನಂತರದಲ್ಲಿ ಕಾಲಸೂಚಕ ಪ್ರತ್ಯಯವೂ ಆಖ್ಯಾತ ಪ್ರತ್ಯಯವೂ ಸೇರಿ ಕರ್ಮಣಿ ಪ್ರಯೋಗವಾಗುವುದು.\par
